\chapter{Saúde e autocuidado mental no mundo digital}

% TODO: Move these refs into bibtex

No quesito de saúde mental, a bipolaridade ou pelo seu nome científico,
Transtorno Afetivo Bipolar (TAB), é uma das doenças mentais mais comuns e
responsável por uma gama muito grande de disfunções mentais. De acordo com
o estudo conduzido por \citet{Merikangas2011}, no âmbito do \textit{World Mental
Health Survey Initiative} da Organização Mundial da Saúde, com mais de 61 mil
adultos entrevistados em 11 países, a prevalência ao longo da vida do espectro
bipolar é de 2,4\% da população mundial, e a prevalência em um período de 12
meses é de 1,5\%. No contexto do Brasil, segundo a Associação Brasileira de
Transtorno Bipolar (ABTB), capítulo brasileiro da \textit{International Society for
Bipolar Disorders} (ISBD) \citep{abtb}, a bipolaridade afeta cerca de
4\% da população, o que representaria uns 7,6 milhões de
pessoas. Esse número, embora soe pouco expressivo, levando em consideração o
contingente de pessoas no país, é preocupante, dado o fato de que a
bipolaridade é (1) complexa de se detectar e (2) muitas vezes, acaba passando
despercebida ou enquadrada dentro de um quadro de depressão unipolar
(depressão).

Embora o autorregistro seja amplamente recomendado no acompanhamento
de pessoas com TAB, dado seu potencial de validade clínica e de efeito
sobre desfechos relevantes~\cite{Faurholt:2016}, sua realização de
maneira contínua representa um desafio.
Em muitos casos, o paciente depende da própria memória para reconstruir
retrospectivamente seu estado emocional, a qualidade do sono,
acontecimentos cotidianos e outros fatores que podem ter influenciado
seu humor. Esse processo está sujeito, naturalmente, a vieses de
recordação e dificulta a identificação de padrões temporais, especialmente
quando as informações são coletadas apenas durante consultas clínicas.

Nesse contexto, o \textit{Ecological Momentary Assessment} (EMA)
consolidou-se como uma importante abordagem metodológica para coleta
de dados em saúde mental. Diferentemente de avaliações retrospectivas,
o EMA propõe que informações sejam registradas no ambiente natural do
indivíduo e próximas ao momento em que as experiências ocorrem, reduzindo
vieses de memória e permitindo observar a dinâmica dos estados afetivos
ao longo do tempo \citep{EbnerPriemer2009}. Além de proporcionar maior
validade ecológica aos dados coletados, essa estratégia tem sido
empregada em diversas áreas da pesquisa comportamental, incluindo
transtornos do humor \citep{Collins2003, EbnerPriemer2009}, permitindo
compreender a relação entre fatores contextuais, eventos cotidianos e
alterações do humor.

Nesse contexto, o uso de tecnologias digitais tem se mostrado uma via
promissora para sistematizar esse tipo de acompanhamento. Pesquisas indicam
que pessoas com transtornos do espectro bipolar demonstram alta adesão ao
monitoramento emocional quando este é feito de forma acessível, visual e
integrada ao seu cotidiano, com estudos reportando taxas de conclusão acima
de 90\% em protocolos de avaliação momentânea
via \textit{smartphone} \citep{Ryan2020} \footnote{
  Vale citar porém, que esse artigo tem limitações de generalização:
  os resultados são limitados a um grupo de pessoas que possuiam um maior
  contato com tecnologia.
}.

Entre as formas mais comuns de implementação do EMA encontra-se o
autorregistro contínuo (\textit{journaling}), no qual o indivíduo
registra sistematicamente aspectos relevantes do cotidiano, como
humor, qualidade do sono, eventos significativos, sintomas percebidos
e outras informações contextuais. Com o acúmulo desses registros,
torna-se possível identificar padrões de comportamento, reconhecer fatores
associados às alterações de humor e favorecer o desenvolvimento da
autoconsciência sobre a própria condição, constituindo uma ferramenta
complementar ao processo de autoconhecimento e autogestão do
comportamento \citep{Melton2013}.

A ampla disseminação dos dispositivos móveis ampliou significativamente
as possibilidades de utilização dessas estratégias fora do ambiente
clínico. Smartphones acompanham o indivíduo durante praticamente toda
sua rotina, permitindo que registros sejam realizados no momento em que
as experiências ocorrem e reduzindo o esforço necessário para manutenção
do acompanhamento contínuo. Como consequência, diversas aplicações
voltadas ao monitoramento da saúde mental passaram a incorporar princípios
do EMA em seus fluxos de utilização, oferecendo registros estruturados
e visualizações históricas capazes de auxiliar usuários e profissionais
de saúde \citep{Faurholt:2016}.

Apesar da crescente disponibilidade de aplicações voltadas ao acompanhamento
do humor, estudos recentes que avaliaram especificamente os aplicativos mais
utilizados por pessoas com TAB para monitoramento de humor e sono apontam
limitações importantes quanto à integração entre esses domínios, ao controle
do usuário sobre seus próprios dados e à acessibilidade das políticas de
privacidade, muitas vezes redigidas em nível de leitura incompatível com o
público geral \citep{Morton2022}.

Em paralelo, análises mais amplas do ecossistema de aplicativos de
saúde mental também identificam riscos recorrentes de privacidade,
como permissões desnecessárias e vazamento de dados sensíveis \citep{Iwaya:22}.
Essas limitações evidenciam a necessidade de soluções que conciliem
boas práticas de engenharia de software, proteção de dados pessoais e
recursos que apoiem a identificação de padrões emocionais ao longo do tempo.

Diante desse contexto, este trabalho propõe o desenvolvimento de uma
aplicação móvel fundamentada nos princípios do
\textit{Ecological Momentary Assessment} voltada ao acompanhamento
cotidiano de pessoas com Transtorno Afetivo Bipolar. O sistema busca
integrar registros estruturados de humor, qualidade do sono, gatilhos
e ações de cuidado em uma única plataforma, produzindo indicadores que
auxiliem o usuário na visualização de padrões e no desenvolvimento da
consciência sobre seu próprio estado afetivo. Ressalta-se que a aplicação
não possui finalidade diagnóstica nem substitui o acompanhamento realizado
por profissionais de saúde, sendo concebida como uma ferramenta de apoio
ao autorregistro e ao autocuidado.
